\begin{theorem}
    \label{carac:riemann}
    Sea $(x_n)_{n\geq1}$ una sucesión de números reales, entonces son equivalentes:
    \begin{enumerate}
        \item $(x_n)$ es equidistribuida módulo $1$.
        \item Para toda función $f\in \mathcal{R}(\mathbb{T})$, Riemann integrable, se cumple que
              \[
                  \lim_{N \to \infty} \frac{1}{N} \sum_{n=1}^N f(x_n) = \int_0^1 f(x)\ dx.
              \]
    \end{enumerate}
\end{theorem}

\begin{proof}
    $1 \Longrightarrow 2$ \, Sea $(x_n)$ una sucesión equidistribuida módulo 1.
    
    Consideremos la partición $0=a_0 < a_1 < \ldots < a_k=1$ de $\mathbb{T}$. Sea $f$ una función 
    escalonada tal que $f(x) = \sum_{i=1}^k c_i \mathbf{1}_{[a_{i-1}, a_i)}(x)$ con $c_i \in \mathbb{R}$
    para cada $i=1,\ldots,k$.

    Entonces,
    \[
        \frac{1}{N} \sum_{n=1}^N f(x_n) = 
        \frac{1}{N} \sum_{n=1}^N \sum_{i=1}^k c_i \mathbf{1}_{[a_{i-1}, a_i)}(x_n) = 
        \sum_{i=1}^k c_i \left(\frac{1}{N} \sum_{n=1}^N \mathbf{1}_{[a_{i-1}, a_i)}(x_n)\right)
    \]
    como $(x_n)$ es equidistribuida, por la observación~\ref{obs:equidistribucion}, 
    \begin{align*}
        \lim_{N \to \infty} \sum_{i=1}^k c_i \left(\frac{1}{N} \sum_{n=1}^N \mathbf{1}_{[a_{i-1}, a_i)}(x_n)\right)
        &= \sum_{i=1}^k c_i \left(\lim_{N \to \infty} \frac{1}{N} \sum_{n=1}^N \mathbf{1}_{[a_{i-1}, a_i)}(x_n)\right) \\
        &= \sum_{i=1}^k c_i \int_0^1 \mathbf{1}_{[a_{i-1}, a_i)}(x)\ dx \\
        &= \int_0^1 \sum_{i=1}^k c_i \mathbf{1}_{[a_{i-1}, a_i)}(x)\ dx 
        =\int_0^1 f(x)\ dx.
    \end{align*}
    
    Ahora, sea $f \in \mathcal{R}(\mathbb{T})$. Como $f$ es integrable Riemann, 
    para todo $\varepsilon > 0$ existe una partición $P=\{0=a_0<\cdots<a_k=1\}$ tal que, 
    si definimos
    \[
        L(x)=\sum_{i=1}^k m_i \mathbf{1}_{[a_{i-1},a_i)}(x) \qquad
        U(x)=\sum_{i=1}^k M_i \mathbf{1}_{[a_{i-1},a_i)}(x)
    \]
    donde $m_i=\inf_{x\in[a_{i-1},a_i)} f(x)$ y $M_i=\sup_{x\in[a_{i-1},a_i)} f(x)$
    se tiene
    \[
        L(x)\le f(x)\le U(x), \qquad \int_0^1 U(x) \, dx - \int_0^1 L(x) \, dx \leq \varepsilon.
    \]

    Sea $S_N = \sum_{n=1}^N f(x_n)$, tenemos,
    \[
        \frac{1}{N} \sum_{n=1}^N L(x_n) \leq \frac{S_N}{N} \leq \frac{1}{N} \sum_{n=1}^N U(x_n)
    \]
    como $L$ y $U$ son funciones escalonadas, al tomar límites,
    \[
        \int_0^1 f(x) \, dx - \varepsilon 
        \leq \int_0^1 L(x) \, dx 
        \leq \liminf_{N \to \infty} \frac{S_N}{N} 
        \leq \limsup_{N \to \infty} \frac{S_N}{N} 
        \leq \int_0^1 U(x) \, dx 
        \leq \int_0^1 f(x) \, dx + \varepsilon
    \]
    como $\varepsilon > 0$ es arbitrario, se concluye que
    \[
        \lim_{N \to \infty} \frac{1}{N} \sum_{n=1}^N f(x_n) = \int_0^1 f(x) \, dx.
    \]

    $2 \Longrightarrow 1$ \, Sea un intervalo $(a,b) \subset \mathbb{T}$ y sea $f(x) = \mathbf{1}_{(a,b)}(x)$, 
    como $f$ es Riemann integrable, entonces,
    \[
        \lim_{N \to \infty} \frac{1}{N} \sum_{n=1}^N \mathbf{1}_{(a,b)}(x_n) = \int_0^1 \mathbf{1}_{(a,b)}(x)\ dx = b-a
    \]

    lo que demuestra que $(x_n)$ es equidistribuida módulo $1$.
\end{proof}

\begin{corollary}
    \label{carac:continua}
    Sea $(x_n)_{n\geq1}$ una sucesión de números reales, entonces son equivalentes:
    \begin{enumerate}
        \item $(x_n)$ es equidistribuida módulo $1$.
        \item Para toda función $f\in C(\mathbb{T})$, se cumple que
              \[
                  \lim_{N \to \infty} \frac{1}{N} \sum_{n=1}^N f(x_n) = \int_0^1 f(x)\ dx.
              \]
    \end{enumerate}
\end{corollary}